\section{Cosmic Government, Action upon Men, Mission, Guardianship, and Demonic Assault}

``Wherever they are sent, they run within him,'' St.~Gregory the Great preached in his \emph{Homily~34 on the Gospels} (\S13): \emph{quolibet missi veniant, intra ipsum currunt}. Because God is uncircumscribed, an angel dispatched to touch a prophet's lips, walk beside Tobias on the road to Rages, or watch over an infant in a cradle never steps outside the divine presence. In Questions 110 through 114 of the \emph{Prima Pars}, St.~Thomas Aquinas traces the entire outward arc of angelic action---how pure spirits govern bodies, enlighten human minds without coercing free will, go forth on divine embassies, guard every human soul along the road of earthly life, and restrain the assaults of the fallen spirits within the bounds of providence.

\subsection{Angelic Government of Corporeal Creation (\ST{I, q.~110})}

All corporeal creatures stand under angelic governance (\ST{I, q.~110, a.~1}). Quoting St.~Augustine (\emph{De Trinitate}~III.4)---``all bodies are ruled by the rational spirit of life''---and St.~Gregory (\emph{Dialogues}~IV.6)---``in this visible world nothing takes place without the agency of the invisible creature''---Aquinas invokes a universal law of order: every particular power bound by matter to ``here'' and ``now'' is governed by a more universal, immaterial power, just as a bailiff serves beneath a king. Bodies act only insofar as they are moved (ad~1). While Aristotle (\emph{Metaphysics}~XI.8) confined separate substances to moving the celestial spheres because he recognized only physical regularity below the moon, Christian faith recognizes providential and supernatural works among earthly bodies, requiring immediate angelic presidency over sublunary creation as well (ad~2). Surveying Plato's separate species-forms, Aristotle's celestial movers, and Avicenna's single active intelligence, Aquinas sides with the Fathers---Augustine (\emph{Eighty-Three Diverse Questions}, q.~79), St.~John of Damascus (\emph{De fide orthodoxa}~II.4), and Origen (\emph{Homily~14 on Numbers}, commenting on Num 22:23)---that divine wisdom sets distinct angels over diverse corporeal domains without multiplying the nine general orders, grouping these cosmic ministers under the Virtues in Gregory's classification (ad~3).

Corporeal matter never obeys an angel's bare nod so as to receive substantial forms directly from angelic thought (\ST{I, q.~110, a.~2}). Against Platonist and Avicennian doctrines of a separate ``giver of forms,'' Aquinas follows St.~Augustine (\emph{De Trinitate}~III.8)---``It is not to be thought, that this visible matter obeys these rebel angels; for it obeys God alone''---and Aristotle (\emph{Metaphysics}~VII.8): what comes into being is not a disembodied form, but a material composite, which demands either a composite bodily agent or God alone, who holds both form and matter in his creative power. Unlike the human soul united as substantial form to its own body (ad~1), an angel alters matter solely by applying natural corporeal agents and hidden material ``seeds'' (\emph{semina}) implanted in the cosmos by God (ad~2--3; Augustine, \emph{De Trinitate}~III.8--9; \emph{De malo} q.~16, a.~9), much as a cook regulates fire to prepare a dish the fire could not produce alone.

The master lever of angelic physics is local motion: bodies obey angels immediately as regards change of place (\ST{I, q.~110, a.~3}; \emph{De malo} q.~16, a.~10). By Pseudo-Dionysius's principle of continuity (\emph{DN}~VII: ``Divine wisdom has joined the ends of the first to the principles of the second'') and Aristotle's proof (\emph{Physics}~VIII.7) that local motion is the first and most perfect bodily movement---altering only extrinsic place rather than intrinsic substance---corporeal nature is naturally apt to be moved locally by spiritual substances. By locally assembling natural agents and seeds, angels initiate qualitative and substantial changes as a smith applies fire to iron (ad~1--2), unrestricted to any single body as the human soul is (ad~3). In \emph{De malo} (q.~16, a.~10, ad~4 and ad~8), Aquinas adds that while demons retain natural power to move sublunary bodies locally, they cannot move the heavenly spheres or shift the entire earth at once.

True miracles belong to God alone (\ST{I, q.~110, a.~4}; Ps 135:4 [Heb 136:4]: ``Who alone doth great wonders''). A miracle properly so called occurs outside the order of the \emph{entire} created nature, not merely outside one particular nature, as when a man throws a stone upward. Because every angelic act stays within the created order, angels ``work miracles'' (St.~Gregory, \emph{Homily~34}, \S10) only by obtaining them from God through intercession or by serving instrumentally, as in gathering human dust at the general resurrection (ad~1). Demonic marvels astonish us merely because their natural speed and use of corporeal seeds outstrip human knowledge; in St.~Augustine's distinction (\emph{Eighty-Three Diverse Questions}, q.~79), magicians act by ``private contracts'' within the cosmos, whereas God's genuine miracles operate by the ``public justice'' of the universal Ruler (ad~2--4).

\angeldossier{\ST{I, q.~110, aa.~1--4} (How angels act on bodies)}{Common teaching (angelic governance of corporeal creation through local motion and natural causes; impossibility of directly informing matter or working true miracles by created power).}{Ps 135:4 [Heb 136:4]; Num 22:23; St.~Augustine, \emph{De Trinitate}~III.4, 8--9, \emph{De div. quaest. LXXXIII} q.~79; St.~Gregory the Great, \emph{Dialogues}~IV.6, \emph{Hom. in Ev.}~34.10; Pseudo-Dionysius, \emph{DN}~VII; St.~John of Damascus, \emph{De fide orthodoxa}~II.4; St.~Thomas Aquinas, \emph{De malo} q.~16, aa.~9--10.}{Platonist cosmology and Avicenna attributed direct formal causation in matter to separate intelligences (a.~2); Aristotle (\emph{Metaphysics}~XI.8) confined separate movers to the celestial spheres (a.~1, ad~2).}

\subsection{Angelic Action upon Human Intellect, Will, Imagination, and Senses (\ST{I, q.~111})}

Divine providence extends hierarchical illumination downward from angels to human beings (\ST{I, q.~111, a.~1}; \CH{IV}). Because human reason in this life understands by turning to sensible images, an angel enlightens a man in two ways: by strengthening his natural intellectual light and by presenting intelligible truth clothed in imaginative phantasms, fulfilling Pseudo-Dionysius's axiom (\CH{I}): ``It is impossible for the divine ray to shine on us, otherwise than shrouded by the variety of the sacred veils.'' God alone infuses the interior habit of faith that moves the will to assent (Eph 2:8; Augustine: ``No one believes except willingly''), while angels ministerially propose truths of faith and practical conduct (ad~1), sharpen reason's ascent from creatures to God (Rom 1:19, ad~2), and enlighten minds that recognize the disclosed truth without necessarily perceiving the angel who brought it (ad~3). In \emph{De malo} (q.~16, a.~12), Aquinas notes that demons share natural angelic power over phantasms yet never truly enlighten human judgment for truth; quoting St.~Augustine (\emph{Eighty-Three Diverse Questions}), they instead cloud the avenues of understanding to provoke pride and error.

The human will, like the angelic will, remains an inviolable sanctuary open directly to God alone (\ST{I, q.~111, a.~2}; Prov 21:1: ``As the divisions of waters, so the heart of the king is in the hand of the Lord: whithersoever he will, he shall turn it''). No created spirit can enter and turn a human will from within. From without, an angel moves human choice solely in two non-coercive ways: by persuading the intellect with an apprehended good, or by stirring bodily passions such as anger or concupiscence in the sensitive appetite, leaving the will completely free to consent or refuse. Reconciling St.~Bede (\emph{Super Matthaeum}~15:11: the devil does not send wicked thoughts, ``but kindles them'') with St.~John of Damascus (\emph{De fide orthodoxa}~II.4: demons ``put'' passions into men), Aquinas shows that demons cannot cause thoughts directly from within since thought obeys the will, yet incite them inwardly through roused passion and suggestion, whereas good thoughts trace principally to God through angelic ministry (ad~1--2). Because human understanding requires phantasms while the will can follow reason against passion, angels exercise greater sway over imagination than over choice (ad~3).

Both good and bad angels can move human imagination by natural power (\ST{I, q.~111, a.~3}), as seen in the dreams guiding St.~Joseph (Matt 1:20; 2:13--19; \emph{De malo} q.~16, a.~11). They never create brand-new sensory species from nothing---an angel ``cannot make a man born blind imagine color'' (ad~1--2)---but employ local motion over bodily spirits and humors (Aristotle, \emph{De insomniis}~III) to recombine stored sensory impressions during sleep, ecstasy, or waking consciousness. St.~Augustine's phrase ``commingling of spirits'' (\emph{De Genesi ad litteram}~XII.12) signifies not a fusion of essences, but causal action on human imagery (ad~3); and if unexplained images lead a receiver astray, the error stems from human weakness rather than any deceit in a good angel (ad~4). In \emph{De malo} (q.~16, a.~11, ad~10), Aquinas adds that demons could naturally agitate phantasms violently enough to block reasoning altogether, were they not held back by divine prohibition.

By the same mastery of local motion, angels can affect the external human senses (\ST{I, q.~111, a.~4}): from without, by presenting a real object or an assumed aerial body (\ST{I, q.~51, a.~2}), and from within, by disturbing bodily spirits and humors as choleric bile makes every taste seem bitter. Thus the angels struck the men of Sodom with dazzling blindness (\gk{aorasia}, Gen 19:11) and blinded the Syrian host at Samaria (4 Kgs [2 Kgs] 6:18), acting beyond a particular bodily routine without breaching the total order of created nature (ad~1--3).

\angeldossiersettled{\ST{I, q.~111, aa.~1--4} (The action of the angels on man)}{Common teaching (illumination of human intellect through phantasms; immunity of the human will from interior coercion; natural power to stir imagination and external senses through local motion of bodily spirits and external forms).}{Prov 21:1; Gen 19:11; 4 Kgs [2 Kgs] 6:18; Matt 1:20, 2:13--19; Rom 1:19, 10:17; Eph 2:8; Heb 1:7; Pseudo-Dionysius, \CH{I, IV}, \emph{EH}~III; St.~Augustine, \emph{De Genesi ad litteram}~XII.12; St.~Bede, \emph{Super Matthaeum}~15:11; St.~John of Damascus, \emph{De fide orthodoxa}~II.4; St.~Thomas Aquinas, \emph{De malo} q.~16, aa.~11--12.}

\subsection{Assisting and Ministering Spirits (\ST{I, q.~112})}

God sends angels on works of external ministry (\ST{I, q.~112, a.~1}; Ex 23:20: ``Behold I will send my angel, who shall go before thee''). Whereas the omnipresent Son and Holy Spirit are ``sent'' by a new mode of presence in grace or Incarnation (\ST{I, q.~43, a.~1}), a finite angel is sent when by divine command he applies his power anew to a corporeal creature where he was not acting before. Aquinas defines a minister as ``an intelligent instrument,'' whose action starts from God's command and ends in God's glory. Stepping out from the empyrean heaven diminishes an angel's dignity no more than leaving a throne diminishes a king's (ad~2), and outward service never clouds contemplation: unlike humans, whose intense bodily labor distracts the intellect, an angel directs outward action through pure intellectual contemplation itself---confirming St.~Gregory (\emph{Moralia}~II) that ``the angels do not go abroad in such a manner as to lose the delights of inward contemplation'' (ad~3). In serving mankind, these loftier spirits serve chiefly God, and honor those united to God by grace (ad~4; Phil 2:3).

Yet not every angelic choir goes forth on external embassies (\ST{I, q.~112, a.~2}). Siding strictly with Pseudo-Dionysius (\CH{XIII}) against theologians who thought God occasionally dispatches the highest spirits by dispensation, Aquinas determines that the supreme hierarchy performs no external service. Unlike the natural order, which yields to grace in miracles to confirm human faith, the celestial order of grace has no higher order above it, and inferior spirits suffice for every human embassy---including the supreme embassy of the Annunciation, committed to the Archangel Gabriel (St.~Gregory, \emph{Homily~34}, \S9). St.~Paul's question, ``Are they not all ministering spirits?'' (Heb 1:14), encompasses either the \emph{interior} mission whereby higher spirits enlighten lower ones, or specifically the angels who mediated the Old Law (ad~1). When Isaias beheld a Seraph purifying his lips with a live coal (Is 6:6--7), Aquinas adopts Dionysius's twin resolution (\CH{XIII}): either an inferior angel was called ``Seraph'' equivocally from his fiery cleansing, or he acted as executor of the Seraphim's power, as a pope absolves through a delegate (ad~2). St.~Gregory (\emph{Homily~34}, \S\S12--13) cites Dionysius with reverence while retaining greater latitude for higher angels in supreme mysteries, insisting that every dispatched spirit remains ``within him'' (\emph{intra ipsum currunt}).

Accordingly, ``assisting'' carries a broad and a strict sense (\ST{I, q.~112, a.~3}). Broadly, every blessed angel assists even while ministering on earth, since none ever loses direct vision of God's essence (Tob 12:15; Matt 18:10; Gregory, \emph{Homily~34}, \S13; \emph{Moralia}~II). Strictly, like royal courtiers hearing the king's decrees directly rather than provincial governors receiving them by messenger, only the entire first hierarchy ``assists,'' enjoying immediate divine illumination into new providential secrets (\ST{I, q.~112, a.~3, resp. and ad~4}; Gregory, \emph{Moralia}~XVII on Job 25:3). When Satan stood among the sons of God (Job 1:6), he did not assist in glory, but appeared merely by surviving angelic nature (ad~3).

Even within the middle hierarchy, the Dominations are not sent externally (\ST{I, q.~112, a.~4}). Being ``above all subjection'' (\CH{VIII}), they dispose and command like architects who never lay hands on stone (ad~1), leaving the five lower orders---Virtues, Powers, Principalities, Archangels, and Angels---to execute external missions. In Reply~2, Aquinas contrasts two classic readings of Daniel 7:10 (``thousands of thousands ministered to him, and ten thousand times a hundred thousand stood before him''). Reading ``thousands out of thousands'' partitively, St.~Gregory (\emph{Moralia}~XVII) concludes that ministering angels outnumber assisting ones, matching six administering choirs to three assisting. Reading the numbers multiplicatively ($10{,}000 \times 100{,}000$ dwarfing $1{,}000 \times 1{,}000$), Pseudo-Dionysius (\CH{XIV}) concludes that assisting angels vastly outnumber ministering ones, because God multiplies nobler spiritual beings more abundantly than all material multitudes.

\angeldossier{\ST{I, q.~112, aa.~1--4} (The mission of the angels)}{Church teaching (a.~1: divine mission of ministering angels); Common teaching (a.~3: uninterrupted beatific vision during external mission); Thomist position (aa.~2, 4: restriction of external mission strictly to the five lower orders, excluding the first triad and Dominations).}{Ex 23:20; Is 6:6--7; Dan 7:10; Tob 12:15; Job 1:6, 25:3; Heb 1:14; Luke 22:27; Phil 2:3; Pseudo-Dionysius, \CH{VII, VIII, XIII, XIV}; St.~Gregory the Great, \emph{Hom. in Ev.}~34.12--13, \emph{Moralia}~II, XVII; St.~Augustine, \emph{De Trinitate}~IV.20; \emph{CCC}~332--335.}{St.~Gregory the Great (\emph{Hom. in Ev.}~34.12--13) allows higher angels to be sent on great mysteries while running within God, whereas Aquinas follows Pseudo-Dionysius (\CH{XIII}) in holding that the four supreme orders never execute external ministry. On a.~4, Gregory's placement of Virtues in the lowest triad yields a different count of ministering choirs from Daniel 7:10.}

\subsection{The Ministry of the Guardian Angels (\ST{I, q.~113})}

Divine providence deputes angels to guard mankind (\ST{I, q.~113, a.~1}; Ps 90:11 [Heb 91:11]: ``For he hath given his angels charge over thee; to keep thee in all thy ways''). Throughout creation, mutable beings are steadied by invariable superiors; because human thoughts are fearful and moral counsels uncertain amid shifting passions (Wis 9:14), free will and natural law do not suffice on their own (ad~1). God's own sleepless watch (Ps 120:4 [Heb 121:4]) guards us immediately by infusing grace and virtue, and mediately by angelic illumination strengthening prudence (ad~2). When sinners perish, the fault rests in human refusal of invisible angelic guidance, never in angelic negligence (ad~3).

More intimately still, each individual human being has a personal guardian angel (\ST{I, q.~113, a.~2}). Unlike irrational animals, which endure only in their species, every human person possesses an incorruptible rational soul destined for eternity; divine providence therefore relates to each individual man as it relates to an entire species of corruptible creatures, assigning a personal angel to each soul. Aquinas invokes Christ's warning---``their angels in heaven always see the face of my Father who is in heaven'' (Matt 18:10)---together with St.~Jerome's celebrated gloss (preserved in Peter Lombard, \emph{II Sent.} d.~11): ``Great is the dignity of souls, for each one to have an angel deputed to guard it from its birth.'' The East teaches the same conviction: St.~Basil the Great (\emph{Adversus Eunomium}~III.1) declares that beside every believer stands an angel ``as pedagogue and shepherd directing his life,'' while St.~Hilary of Poitiers (\emph{Tractatus in Psalmum~124}) keeps the creaturely hierarchy clear: ``Good indeed is the protection of an angel, but better that of the Lord.'' Whereas one human ruler suffices for external communal order, personal angelic guardianship attends to the hidden interior path of each soul's salvation, with greater offices matching greater providential vocations (ad~1--3; Ecclus 33:11--12).

Personal guardianship over an individual belongs properly to the lowest order of Angels (Ps 90:11; \CH{V, IX}), while higher choirs hold universal care---Principalities or Archangels (such as Michael, ``one of the chief princes,'' Dan 10:13) over nations, Virtues over bodies, and Powers restraining demons (\ST{I, q.~113, a.~3}). Where St.~John Chrysostom (\emph{Homily~59 on Matthew}, \S4) reads Matthew 18:10 as speaking ``not of angels only, but rather of angels that are greater than others,'' Aquinas accommodates him as referring either to the loftiest spirits within the lowest order (\CH{X}) or to higher guardians assigned to those elected to extraordinary glory (ad~1); moreover, lower guardians curb demons and perform wonders by sharing in higher choirs' power (ad~2--3).

Every earthly wayfarer receives a guardian angel along the snare-strewn path of life (Ps 141:4 [Heb 142:4]; \ST{I, q.~113, a.~4}). Adam had a guardian in Eden because external demonic snares threatened even unfallen innocence (ad~2); unbelievers, the reprobate, and even Antichrist (2 Thess 2:9) retain their guardians just as they retain natural reason, warding off countless harms they would otherwise commit (ad~3). At death guardianship concludes: the saved reign beside an angel in heaven, the lost suffer beneath a demon. Christ alone stands outside this rule (ad~1): guided immediately by the Word and already a \emph{comprehensor} in his human soul, he had no superior guardian angel, though in his passible humanity inferior angels came and ministered to him (Matt 4:11).

When personal guardianship begins was an open inquiry among the Fathers (\ST{I, q.~113, a.~5}). In his \emph{Commentary on Matthew} (XIII.26--29), Origen asked whether an angel takes charge of a child at birth or at baptism---and even floated a speculative hypothesis, rejected by Catholic faith, that a wicked angel assigned at birth might later convert. Siding with St.~Jerome, Aquinas assigns personal guardianship from \emph{birth}: supernatural sacraments begin at baptism, whereas angelic guardianship is a providential benefit bestowed upon human rational nature itself as soon as the infant becomes an individual separate from the mother (ad~1--2), guarding even pre-rational infants from bodily and demonic harm. While the child remains in the womb, Aquinas judges with measured probability that the mother's guardian angel protects the unborn child (ad~3).

A guardian angel never wholly forsakes his charge during mortal life (\ST{I, q.~113, a.~6}; 1 Pet 5:8). Scriptural words of angelic withdrawal (Jer 51:9; Is 5:5) signify particular permissions of tribulation or sin under divine justice, never total abandonment. Reconciling St.~John of Damascus's remark that angels on earth are not locally in heaven (\emph{De fide orthodoxa}~II.3), Aquinas replies that even when locally in heaven, a guardian angel knows his charge's state and can be locally present in an instant (ad~3).

Nor do blessed angels suffer painful grief when their charges sin or perish (\ST{I, q.~113, a.~7}; Apoc 21:4: ``nor mourning, nor crying, nor sorrow''). Because sorrow arises from what happens against one's will (St.~Augustine, \emph{City of God}~XIV.15), and blessed angels cleave wholly to divine justice---willing absolutely whatever God effects or permits, as a sailor wills to jettison cargo in a storm (Aristotle, \emph{Nicomachean Ethics}~III.1)---nothing occurs simply contrary to their will. Isaias's weeping ``angels of peace'' (Is 33:7) denote literally Ezechias's envoys, allegorically preachers mourning sin, and anagogically the blessed angels' antecedent desire for all men's salvation (ad~1); angelic joy over repentance (Luke 15:7) rests on the fulfillment of providence (ad~3), and angels stand at judgment not as defendants, but as witnesses (ad~4).

Finally, Aquinas resolves the seeming angelic conflict in Daniel 10:13: ``The prince of the kingdom of the Persians resisted me one and twenty days'' (\ST{I, q.~113, a.~8}). Among good angels, discord of will is impossible, for God ``maketh peace in his high places'' (Job 25:2). Aquinas records two patristic readings: he cites St.~Jerome for identifying the prince of Persia with an evil spirit hindering Israel's release by tempting captives into sin (while in Jerome's own \emph{Commentary on Daniel}~10:13, Jerome identifies him as the angel entrusted with Persia under Deut 32:8 LXX acting on behalf of his province), and places beside it St.~Gregory's reading (\emph{Moralia}~XVII) that Persia's prince was a good guardian angel. On Gregory's view, good angels ``resist'' one another solely in intellectual consultation, presenting the rival merits of nations before divine wisdom until God reveals his decree---whereupon every angelic will concurs in perfect unity.

\angeldossier{\ST{I, q.~113, aa.~1--8} (Of the guardianship of the good angels)}{Church teaching (aa.~1--2: providential angelic guardianship and personal guardian angel for every human person); Common teaching (aa.~4, 6--8: universal guardianship of wayfarers, Christ's exemption as comprehensor, constancy of guardianship, absence of painful grief in glory, concord of angelic wills); Thomist position (a.~3: assignment of personal guardianship to the lowest order of Angels); Disputed (a.~5: commencement of personal guardianship at birth versus baptism).}{Ps 90:11--12 [Heb 91:11--12]; Ps 120:4 [Heb 121:4]; Ps 141:4 [Heb 142:4]; Wis 9:14; Is 5:5, 33:7; Jer 51:9; Dan 10:13; Matt 4:11, 18:10; Luke 15:7; Heb 1:14; 1 Pet 5:8; Apoc 21:4; St.~Basil, \emph{Adversus Eunomium}~III.1; St.~Hilary, \emph{In Ps.}~124; St.~John Chrysostom, \emph{Hom. in Matt.}~59.4; St.~Jerome, \emph{In Matthaeum}~18:10, \emph{In Danielem}~10:13; St.~Augustine, \emph{De civitate Dei}~XIV.15; St.~Gregory the Great, \emph{Hom. in Ev.}~34, \emph{Moralia}~XVII; St.~John of Damascus, \emph{De fide orthodoxa}~II.3; \emph{Roman Catechism} (1566), IV, Preface; \emph{CCC}~336.}{Origen (\emph{In Matthaeum}~XIII.26--29) leaves open whether personal guardianship begins at birth or at baptism (a.~5), whereas St.~Jerome and Aquinas assign it from birth; St.~John Chrysostom (\emph{Hom. in Matt.}~59.4) speaks of greater angels guarding believers without restricting personal guardianship to the lowest choir (a.~3). On Dan 10:13 (a.~8), Aquinas reports Jerome's demonic-prince reading alongside St.~Gregory's good-guardian reading (\emph{Moralia}~XVII), while Jerome's own \emph{In Danielem}~10:13 identifies Persia's presiding angel.}

\subsection{Demonic Assault Within Divine Providence (\ST{I, q.~114})}

The treatise closes with the warfare of the pilgrim Church (\ST{I, q.~114, a.~1}; Eph 6:12; Council of Trent, Sess.~6, ch.~13). Two distinct sources meet in every demonic assault: demonic envy and pride, which parody divine mission by assigning evil spirits to assail men, and divine providence, which turns evil to the testing and crown of the elect. God never sends demons to tempt to sin, only permits temptation by just judgment, though he sometimes dispatches them as instruments of temporal punishment---like the lying spirit sent against Achab (3 Kgs [1 Kgs] 22:20)---where God acts in justice while the demon acts in hatred (ad~1). Nor is the combat unequal: human frailty is armed chiefly with divine grace and secondarily with angelic guardianship, as Eliseus told his servant: ``Fear not: for there are more with us than with them'' (4 Kgs [2 Kgs] 6:16; ad~2--3).

To tempt is properly to make trial in order to know (\ST{I, q.~114, a.~2}). God tests men to manifest their fidelity to others (Gen 22:1; Deut 13:3); the world and the flesh tempt as material occasions; only the devil has tempting-to-ruin as his proper office (1 Thess 3:5). Because God alone weighs the heart (Prov 16:2), demons test men to discover which hidden inclination exposes a soul to sin (ad~2), inclining choice only through imagination and passion without ever forcing the will (ad~3).

Not every human sin springs directly from demonic instigation (\ST{I, q.~114, a.~3}). Indirectly and occasionally, the devil caused all our sins by provoking the Fall that wounded human nature---saving the strong sayings of Pseudo-Dionysius (\emph{DN}~IV) and St.~John of Damascus (\emph{De fide orthodoxa}~II.4). Directly, however, citing Gennadius (\emph{De ecclesiasticis dogmatibus}~XLIX) and Origen (\emph{Peri Archon}~III), Aquinas observes that even if no devil existed, disordered bodily appetites and uncurbed free will would still produce sin, making sinners children of the devil by imitation (John 8:44, ad~2). Yet good and evil are asymmetric: man can fall by himself, whereas he cannot merit eternal life without divine grace mediated through good angels, who cooperate in every meritorious work (ad~3).

Demons can never deceive mankind with genuine miracles in the strict sense (\ST{I, q.~114, a.~4}; \ST{I, q.~110, a.~4}). What Scripture calls ``lying wonders'' (2 Thess 2:9) are either real physical prodigies accelerated through hidden corporeal seeds---such as Pharaoh's magicians' serpents and frogs, or the fire and tempest striking Job's household (St.~Augustine, \emph{City of God}~XX.19; \emph{De Trinitate}~III.8--9; \emph{De malo} q.~16, a.~9)---or sensory and imaginative illusions whenever something exceeds nature, such as changing men into beasts or raising the dead (Augustine, \emph{City of God}~XVIII.18; ad~1--2). True miracles and demonic magic stand apart in both end and right: God's glory under public justice versus pride under private compact (ad~3; Augustine, \emph{Eighty-Three Diverse Questions}, q.~79).

Finally, a demon overcome by a Christian is not permanently disabled from returning, yet withdraws for a definite time (\ST{I, q.~114, a.~5}; Matt 4:11; Luke 4:13: ``departed from him for a time,'' weighed against Matt 12:44). Aquinas cites both God's mercy measuring assault to human frailty (\emph{Opus Imperfectum in Matthaeum}, Hom.~5) and Satan's own astuteness shrinking from repeated defeat (St.~Ambrose, \emph{Exposition of Luke}~4:13, cited in \ST{I, q.~114, a.~5}). Liturgical rites of exorcism, canonical discipline, and pastoral discernment belong to the separate study on Catholic exorcism.

\angeldossiersettled{\ST{I, q.~114, aa.~1--5} (Of the assaults of the demons)}{Church teaching (a.~1: real demonic assault permitted by Providence for the exercise and crown of virtue); Common teaching (aa.~2--5: temptation to ruin as proper office of the devil, distinction between direct demonic instigation and human concupiscence, deceptive prodigies versus true miracles, temporary restraint of defeated demons).}{Gen 22:1; Deut 13:3; Job 1; 3 Kgs [1 Kgs] 22:20; 4 Kgs [2 Kgs] 6:16; Prov 16:2; Matt 4:11, 12:44; Luke 4:13; Mark 16:20; John 8:44; Eph 6:12; 1 Thess 3:5; 2 Thess 2:9; Pseudo-Dionysius, \emph{DN}~IV; St.~Augustine, \emph{De civitate Dei}~XVIII.18, XX.19, \emph{De Trinitate}~III.8--9, \emph{De div. quaest. LXXXIII} q.~79; Gennadius, \emph{De ecclesiasticis dogmatibus}~XLIX; St.~Ambrose, \emph{Expositio in Lucam}~4:13 (cited in \ST{I, q.~114, a.~5}); St.~John of Damascus, \emph{De fide orthodoxa}~II.4; Council of Trent, Sess.~6, ch.~13; CDF, \emph{Christian Faith and Demonology} (1975); \emph{CCC}~395.}
